\appendix
\section{Шаблон журнала эксперимента}

Для воспроизводимости каждый запуск должен сопровождаться записью параметров. Минимальная форма приведена в табл.~\ref{tab:run_log}.

\begin{table}[H]
\centering
\caption{Шаблон журнала одного запуска}
\label{tab:run_log}
\begin{tabularx}{\textwidth}{lY}
\toprule
Поле & Значение \\
\midrule
Идентификатор запуска & \todo{строка} \\
Дата запуска & \todo{дата} \\
Корпус & \todo{название} \\
Словарь & \todo{описание} \\
Основа модели & \todo{название малой модели} \\
Максимальное число шагов & \todo{число} \\
Семейство распределения времени & \todo{название} \\
Правило остановки & \todo{название} \\
Порог правила & \todo{число} \\
Коэффициент стоимости & \todo{число} \\
Начальное зерно & \todo{число} \\
Примечания & \todo{кратко} \\
\bottomrule
\end{tabularx}
\end{table}

\section{Минимальный список проверок перед защитой}

Перед включением численных результатов в основной текст нужно проверить четыре вещи. Во-первых, все методы должны использовать один и тот же словарь и одно и то же разбиение данных. Во-вторых, пороги правил остановки должны подбираться только на проверочной выборке. В-третьих, сравнение методов должно проводиться не только при одном пороге, но и при одинаковом среднем бюджете. В-четвертых, для всех таблиц должны быть указаны средние значения по нескольким запускам и разброс результатов.

\end{document}
